\documentclass[17pt,a4paper,landscape]{extarticle}
\usepackage[margin=2.35cm,top=0.12cm,left=1.05cm]{geometry}
\usepackage{amsmath}
\usepackage{amssymb}
\usepackage{tasks}
\usepackage{tikz}
\usepackage{xcolor}
\usepackage[sfdefault,lf]{FiraSans}
\renewcommand{\familydefault}{\sfdefault}
\setlength{\parindent}{0pt}
\pagestyle{empty}
\settasks{label=\textbf{\arabic*.}, label-width=2.2em, item-indent=2.95em, column-sep=2.1cm, after-item-skip=4.7em}
\renewcommand{\arraystretch}{1.15}
\everymath{\displaystyle}

\begin{document}
\sffamily
\boldmath

\boldmath

{\LARGE \textbf{Area of trapezium — Foundation}}
\begin{tasks}(3)
\task {Identify and circle the two parallel sides in the drawn trapezium.}
\task {Calculate the area of a trapezium with parallel sides of 7 m and 18 m, and a perpendicular height of 6 m.}
\task {Explain why the area formula for a trapezium can be thought of as "the average of the parallel sides multiplied by the height."}
\task {Identify and circle the two parallel sides in the drawn trapezium.}
\task {Calculate the area of a trapezium with parallel sides of 3 m and 15 m, and a perpendicular height of 12 m.}
\task {Explain why the area formula for a trapezium can be thought of as "the average of the parallel sides multiplied by the height."}
\task {Identify and circle the two parallel sides in the drawn trapezium.}
\task {Calculate the area of a trapezium with parallel sides of 1 m and 13 m, and a perpendicular height of 12 m.}
\task {Explain why the area formula for a trapezium can be thought of as "the average of the parallel sides multiplied by the height."}
\task {Identify and circle the two parallel sides in the drawn trapezium.}
\task {Calculate the area of a trapezium with parallel sides of 20 m and 16 m, and a perpendicular height of 10 m.}
\task {Explain why the area formula for a trapezium can be thought of as "the average of the parallel sides multiplied by the height."}
\task {Identify and circle the two parallel sides in the drawn trapezium.}
\task {Calculate the area of a trapezium with parallel sides of 15 m and 5 m, and a perpendicular height of 10 m.}
\end{tasks}

\end{document}
